\begin{problem}
    En la figura adjunta se explican los términos “DAP” y “altura comercial” que se utilizan en la industria de producción forestal.

    \begin{center}
    (imagen)
    \end{center}

    El volumen de madera que se obtiene de un cedro se calcula con la expresión
    \[
    \frac{3}{4}\pi \left(\frac{\text{DAP}}{2}\right)^2 \cdot \text{altura comercial}.
    \]

    ¿Cuál es el volumen de madera que se obtendrá de un cedro de altura comercial de \SI{6}{m} y DAP de \SI{0.4}{m}?

    \begin{enumerate}[label=\Alph*.]
        \item $\SI{0.72}{\pi m^3}$
        \item $\SI{0.18}{\pi m^3}$
        \item $\SI{0.039}{\pi m^3}$
        \item $\SI{0.156}{\pi m^3}$
    \end{enumerate}
\end{problem}
